\section{Discussion}

The frozen analysis \PrimaryInterpretation. This statement is deliberately
stronger than asking whether CFR-WM's point estimate is favorable: H1 required
the lower bootstrap endpoint to clear a 5\% practical margin and the exact
seed-level test to reject that same margin. The distinction matters with eight
top-level realizations, for which apparently stable trajectory-level averages
can coexist with substantial training-seed uncertainty. The result therefore
answers a narrow question about equal-trust context correction on the fixed
benchmark; it is neither a universal ranking of world-model architectures nor
a post hoc selection of the most favorable comparator.

The regime decomposition gives the proposed mechanism a more specific stress test. The
corrupt/shifted estimate was \CorruptImprovement\%, whereas the clean estimate
was \CleanImprovement\%; after hierarchical gating, the full analysis
\RobustnessInterpretation. Context corruption and parameter shift stress
different terms in the error recursion. Noise and missingness weaken the
measurement term in Equation~\eqref{eq:filter-mse}, while parameter shift biases
the nominal pre-update error and the subsequent nominal rollout. A gain of 0.5
cannot adapt between these cases. Consequently, similar aggregate effects can
mask different cell patterns. The three cells with more than 10\% degradation
were all clean, led by cart--pole, while every ID-corrupt cell improved
descriptively. This pattern is consistent with noise suppression within the
bundled corrupt regimes, but the design cannot separate Gaussian noise from
missingness or parameter shift; equal trust can also over-correct already
accurate clean context. This is why the regime effects and all 15 cells are
retained without promoting them to independent confirmatory or causal claims.

The gain sweep further prevents the fixed value from being interpreted as an
optimal filter. The descriptively lowest aggregate NRMSE among the five frozen
gains occurred at \BestGain. This result locates equal trust within a continuum
from pure nominal context propagation (gain 0) to direct measurement
replacement (gain 1). It does not estimate a Bayes-optimal gain, because the
noise, nominal-model error, and missingness processes vary by channel and
trajectory. A learned gain was considered during development, but its advantage
over the matched fixed-gain model missed a prespecified mechanism threshold;
excluding it from confirmation avoids relabeling residual-model capacity as
filter adaptation.

Multi-step training addresses a different part of the recursion. CFR-WM's
relative effect against the one-step-trained variant was
\OneStepImprovement\% (95\% CI [\OneStepRelativeCILow\%,
\OneStepRelativeCIHigh\%]). The context filter changes the initial-error term,
whereas multi-step likelihood exposes the residual transition to its own
forecast distribution and targets accumulated local errors. Their roles are
therefore complementary in the mathematical sense of the bound, even if the
Holm-adjusted sensitivity contrast does not support a population-level
component claim.

Relative to recent world-model and physics-guided work, the value of CFR-WM is
auditability rather than architectural scale. It uses a transparent
measurement rule, explicit missing-channel semantics, known-control free
rollout, and a fixed practical-effect gate. This design trades adaptive
expressivity for a mechanism that can be falsified with approximately
parameter-matched gains.
The archived failure of two earlier candidates is part of that evidence: it
shows that physics, uncertainty, or a learned gate cannot be treated as useful
merely because each is plausible in isolation. Future extensions should adopt
the same separation when adding visual encoders, structured occlusion,
ensembles or conformal calibration, and closed-loop planning
\cite{portela2025_conformal_prediction_for_uncertainty_quantification,brunke2022_safe_learning_in_robotics_from}.

Most importantly, the study distinguishes mathematical, simulated, and
deployment evidence. The filter identity is conditional; the benchmark
establishes measured predictive effects under explicit perturbations; neither
establishes causal mechanism, policy benefit, or hardware safety. Maintaining
those boundaries is necessary for a compact physical-world-model result to be
scientifically useful rather than rhetorically broader than its data.
